% ═══════════════════════════════════════════════════════════════════════
% by1-en.tex — English full paper.  Compile: pdflatex/latexmk (Overleaf OK).
% Chinese counterpart: by1-zh.tex.  Shared references: by1-refs.bib.
% ═══════════════════════════════════════════════════════════════════════
\documentclass[11pt,a4paper]{article}

\usepackage[margin=2.6cm]{geometry}
\usepackage{booktabs}
\usepackage{array}
\usepackage{amsmath}
\usepackage{listings}
\usepackage{microtype}
\usepackage[hidelinks]{hyperref}

\lstset{
  basicstyle=\ttfamily\small,
  columns=fullflexible,
  keepspaces=true,
  frame=single,
  framerule=0.4pt,
  xleftmargin=1.2em,
  xrightmargin=1.2em,
  aboveskip=0.9em,
  belowskip=0.9em
}

\newcommand{\code}[1]{\texttt{#1}}
\newcommand{\byone}{\textsc{by1}}
\newcommand{\ir}{\textsc{by1-ir}}
\newcommand{\q}[1]{\textquotedblleft #1\textquotedblright}

\title{\byone{}: Declarative Descriptions of Large-Model Architectures\\
       with Machine-Verifiable Contracts}
\author{Draft v0.1 --- author list and affiliation to be completed}
\date{October 2026}

\begin{document}
\maketitle

\begin{abstract}
Large language models are described today in several carriers at once: a
configuration file, a set of tensor names, converter code, runtime code, and
prose documentation. The same architectural fact---how many key--value heads a
layer has, which parameters are shared, what a normalization epsilon is---must
be repeated in every carrier, and the carriers drift apart silently. We present
\byone{}, a small declarative language for describing model architectures and
checkpoint contracts, together with a backend-neutral intermediate
representation (\ir{}) and a validation pipeline that grades a description
against artifacts the description did not produce: the publisher's
configuration files and tensor manifests, and independently maintained
runtimes. \byone{} refuses to guess: unknown mechanisms and inconsistent
configurations are errors, and any combination a backend cannot compute is
rejected rather than approximated. Support is reported in four separately
evidenced states---specification check, checkpoint load, numeric run, exported
artifact---so that wiring evidence is never promoted to runtime evidence. On
six published checkpoints (GPT-2, Qwen3-0.6B, Qwen2.5-7B-Instruct, Qwen3-8B,
\code{gpt-oss-20b} in MXFP4, and Qwen3.8-27B with a vision tower and multi-token
prediction), models compiled from \byone{} descriptions match reference
Transformers implementations with a maximum absolute logit difference of
\textbf{0} and identical greedy token sequences; one real training step matches
596{,}049{,}920 gradients and the same number of updated parameters exactly. We
report these as evidence that description, conversion and weight loading are
faithful; because the experiments reuse the reference operators, they say
nothing about kernel-level correctness or speed. We describe the language, the
checking methodology, and the failure modes---including one in which an
incorrect description passed every check---that shaped both.
\end{abstract}

\section{Introduction}

A modern model release arrives as a pile of mutually redundant artifacts. There
is a \code{config.json} with a few dozen hyperparameters; a weight file whose
tensor names encode the architecture implicitly; one or more converters that
map those names into a runtime's conventions; a model class that encodes what
is actually computed; and a README or model card that says, in prose, what the
authors believe they built. Each of these carriers is a partial description of
the same object, and none of them is authoritative over the others.

This redundancy is not a curiosity; it is the dominant source of expensive,
invisible errors. A bias flag that differs between the query projection and the
output projection changes the tensor inventory. A normalization epsilon that
some converter writes as a local constant can differ from the checkpoint's.
Parameter tying---an output head that is literally the same tensor as the input
embedding---is invisible in a shape listing. Vision towers, multi-token
prediction heads (MTP), sliding-window schedules and mixture-of-experts
routing each add carriers of their own. When carriers disagree, nothing
happens loudly: the model loads, runs, produces plausible text, and is subtly
wrong.

\byone{} takes the position that the architecture itself deserves a single,
small, machine-checkable description. Its surface language (files with the
suffix \code{.by1}) declares hyperparameters, layer patterns, mechanisms,
tensor contracts, sharing relations, positional schemes, state contracts and
export rules. The declaration compiles to \ir{}, a backend-neutral intermediate
representation whose legality is defined by a closed set of operator kinds and
a per-kind attribute table. The validation pipeline then grades the
description against \emph{external} artifacts---the publisher's own
configuration and tensor manifest, and independently maintained
implementations---rather than against itself.

Three design commitments do most of the work.

\begin{itemize}
  \item \textbf{Implement it or reject it.} There is no third path. A
  mechanism that no backend can compute must be refused by the checker, not
  silently lowered into something that happens to run. This rule exists
  because the project violated it once (Section~\ref{sec:lessons}) and passed
  every check with a description that computed the wrong model.
  \item \textbf{Evidence is graded, not counted.} \q{Supported} is split into
  four states---specification check, checkpoint load, numeric run, exported
  artifact---each with its own evidence. Random-weight shadow comparisons are
  recorded as wiring evidence and never promoted to runtime evidence. Numbers
  rise only when evidence rises.
  \item \textbf{The specification is generated where possible.} The \ir{}
  reference document is produced from the same tables the validator executes;
  the language specification's examples are fed back into the checker, and
  green examples must pass while red examples must fail with specified
  diagnostics. A hand-written spec that can go stale is treated as worse than
  no spec.
\end{itemize}

We make the following contributions:
\begin{enumerate}
  \item A declarative architecture language and a backend-neutral IR whose
  legality is a closed, table-driven set (Sections~\ref{sec:lang} and~\ref{sec:ir}).
  \item A verification methodology that grounds claims in publisher artifacts
  and external runtimes, with explicit support states and negative-test
  discipline (Section~\ref{sec:verify}).
  \item An implementation with four cross-checked backends---NumPy, PyTorch,
  generated C, generated \code{ggml}---plus a Transformers backend for
  complete published models (Section~\ref{sec:impl}).
  \item An evaluation on six real checkpoints, including an MXFP4
  mixture-of-experts model and a 27B multimodal model with MTP, with exact
  logit parity and a matched real training step (Section~\ref{sec:eval}).
  \item An account of the negative results and process lessons, including the
  false-positive incident and the provenance of the code
  (Section~\ref{sec:lessons}).
\end{enumerate}

\section{Problem and Scope}
\label{sec:problem}

\subsection{Where architecture knowledge lives today}

Table~\ref{tab:carriers} summarizes the carriers and their characteristic
failure. The pattern is uniform: every carrier holds part of the truth, and
the binding between carriers is convention, not contract.

\begin{table}[t]
\centering
\caption{Architecture knowledge and its carriers. \byone{} replaces the
binding column with explicit contracts.}
\label{tab:carriers}
\begin{tabular}{@{}p{3.2cm}p{4.6cm}p{5.4cm}@{}}
\toprule
Carrier & Holds & Typical failure \\
\midrule
Config file & hyperparameters, some flags & fields that no consumer reads;
fields that consumers invent \\
Tensor names/shapes & implicit structure & tied weights look duplicated;
absent biases look present \\
Converter code & naming and layout rules & per-runtime reimplementation;
silent defaulting \\
Model class & what is computed & disagrees with config; undocumented
overrides \\
Prose docs & intent & drifts from all of the above \\
\bottomrule
\end{tabular}
\end{table}

\subsection{Two failure modes we target}

\paragraph{Silent defaulting.} When a converter meets an unknown or missing
field, the path of least resistance is to pick a plausible value. The project
record contains a concrete instance: normalization epsilon was hard-coded in
three separate places before it became a required IR field. The failure mode
is quiet because a wrong epsilon rarely changes token sequences in a short
greedy run; it changes numerics just enough to be deniable.

\paragraph{Name-and-shape agreement is not computational agreement.}
Early in the project, \code{gpt2.by1} compiled to byte-identical IR as
\code{llama-shaped.by1}---a GPT-2 description and a Llama-style description
producing the same intermediate representation. Every backend, the verifier
and the full check suite reported \q{pass}: tensor contract checks verified
names and shapes, not what is computed. Learned positional embeddings had been
lowered into an abstraction that also covered rotary positions. The fix was
not a better check but a changed rule: the IR's positional kind is an enum
that cannot express both, and any mechanism that lacks a real implementation
is rejected. Section~\ref{sec:lessons} returns to this.

\subsection{Goals and non-goals}

The intended users are (i)~model publishers and converters who must turn a
checkpoint into auditable standard artifacts, (ii)~architecture reviewers who
must confirm that documentation, configuration and stored weights agree, and
(iii)~re-implementers who should be able to build a backend from the IR
specification without reading \byone{}'s source.

Equally important is what \byone{} deliberately is \emph{not}: not a general
inference runtime, not a training framework, not a serving system, and not an
author of CUDA kernels. Deployment uses established runtimes
(Transformers~\cite{wolf2020transformers}, \code{llama.cpp}~\cite{llamacpp2023});
quantized artifacts use maintained formats rather than private ones; and
private extensions to a standard format must first demonstrate that the
standard cannot express the case, with unknown extensions rejected on read.

\section{Language Design}
\label{sec:lang}

\subsection{Surface language}

A complete architecture is a short document. Listing~\ref{lst:hello} is the
project's CPU demo in full; it runs on a laptop with NumPy alone and random
weights.

\begin{lstlisting}[caption={The complete demo architecture
\code{models/hello.by1}.}, label={lst:hello}]
model Hello {
  hparams { d_model = 128, vocab = 256 }
  stack main {
    pattern = 4 * Attention(q = 4, kv = 2, head_dim = 32)
    main[:] >> FFN(hidden = 512)
  }
}
\end{lstlisting}

The declaration is intentionally small, but it is not a lossy summary: the
checker resolves it to a full tensor contract (four decoder layers, fused
query/key/value projections with the declared biases, a gated feed-forward
block attached to every layer, and the default positional, normalization and
head rules). Table~\ref{tab:onechange} reports the effect of changing the
single digit in the pattern line.

\begin{table}[t]
\centering
\caption{One number changed in Listing~\ref{lst:hello}. The layer count and
the parameter budget follow the declaration; the output shape does not change
because the vocabulary and sequence length do not.}
\label{tab:onechange}
\begin{tabular}{@{}lrr@{}}
\toprule
Result & \code{4 * Attention(...)} & \code{2 * Attention(...)} \\
\midrule
Decoder layers & 4 & 2 \\
Parameters allocated & 1{,}049{,}728 & 557{,}696 \\
Output shape & $(1, 8, 256)$ & $(1, 8, 256)$ \\
\bottomrule
\end{tabular}
\end{table}

Beyond the skeleton---\code{model}, \code{hparams}, \code{stack},
\code{pattern}, \code{position}, \code{state}, \code{optimizer}---the language
carries two kinds of statement that simple config formats lack.

\paragraph{Mechanisms with declared defaults and invariants.}
A \emph{mechanism} is a named computational unit (\code{Attention},
\code{FFN}, \code{KDA}, \code{MLA}, \code{MoE}, \code{SSM}, and others) with a
parameter table. Mechanism names come from a single registry; writing a name
that the registry does not know is an error, by design.

\paragraph{Pattern scheduling over instances.}
\code{pattern} schedules repetitions and alternations of mechanism instances,
e.g.\ a ratio of three gated-convolution layers to one attention layer, or
alternating sliding-window and full attention. Attribute values on a pattern
element are classified by their downstream effect, because the three classes
deserve different rules:
\begin{enumerate}
  \item \emph{behavioral} attributes (window size, sink, gating mode) change
  how an instance computes but not its tensor inventory;
  \item \emph{structural} attributes (head dimension, key--value head count,
  tying) change shapes and the set of tensors, so an instance that overrides
  one must be re-validated and its tensor contract expanded per instance;
  \item \emph{scheduling} attributes point at another instance's state
  (a layer reading a previous layer's cache).
\end{enumerate}

\subsection{Tensor contracts and sharing}

A description states, per block, which tensors exist, their shapes, dtypes and
sharing relations. Sharing is a first-class relation: input embedding and
output head may be the same object, may be distinct with different biases, and
the relation is checked against the checkpoint's actual tensor inventory
rather than inferred from shapes. The six complete models in the corpus were
selected precisely because they exercise these differences: shared versus
independent output heads, projection bias asymmetries in Qwen2.5-7B-Instruct,
learned versus rotary positions in GPT-2 and Qwen, mixture-of-experts with
sliding-window alternation in \code{gpt-oss-20b}, and a 27B model combining a
text tower, a vision tower and an MTP head.

\subsection{State contracts}

Recurrent and cache-carrying mechanisms declare state explicitly: kind
(key--value cache, recurrent state, convolution history, or declared by the
implementation), whether the state is bounded and by what, its shape and
dtype, and whether it may be reused across requests (none, prefix, or
unknown). Making state a declaration rather than an implementation detail is
what allows a reviewer to compute cache sizes from the description alone, and
it is the basis of the state-size hand-calculation task in the audit procedure
(Section~\ref{sec:verify}).

\section{The Intermediate Representation}
\label{sec:ir}

\ir{} is the contract between the front end and every backend. Its top level
holds global fields (vocabulary size, context length, residual width,
positional kind, normalization kind and epsilon, and so on), a registry of
stacks (text, vision, MTP), the global operator list (embedding, final
normalization, head), and the per-layer list of operators with resolved
attributes and state declarations.

Two rules give the IR its discipline:

\paragraph{Closed sets.} Operator kind is a closed set of thirteen kinds
(\code{Add}, \code{Attention}, \code{Embed}, \code{External}, \code{FFN},
\code{Head}, \code{KDA}, \code{Linear}, \code{MLA}, \code{MoE}, \code{Norm},
\code{Raw}, \code{SSM}); each kind has a table of required and permitted
attributes. An attribute absent from the table is illegal; a required
attribute that is missing is illegal. State dtypes and reuse modes are
likewise enums: a free-form string would make \q{unknown value} an untestable
state, and the project had already been bitten by that.

\paragraph{Generated specification.} The reference document for \ir{} is
emitted from the same tables that the validator executes
(\code{python src/by1ir.py --spec}). The rationale is stated in the document
itself: a hand-written specification goes stale, and a stale specification is
worse than none because it makes the reader believe they know. Readers must
reject IR major versions they do not recognize; version history is
append-only, and existing \code{1.0} files remain valid under \code{1.1}.

\subsection{Error model and versioning}

Diagnostics are one per line with a location and an actionable reason; any
error yields a nonzero exit status. Silent fallback is prohibited: an unknown
mechanism, an unknown dtype, or a contradictory configuration is an error,
never \q{the nearest guess}. The negative-example baseline
(\code{models/selftest.by1}) is pinned at exactly nine errors, so a parser
change that alters the error count must update the baseline explicitly and
explain why. The language specification document doubles as a test suite: its
green examples must pass the checker and its red examples must fail with
diagnostics containing specified substrings.

\section{Verification Methodology}
\label{sec:verify}

\subsection{External graders}

A description is graded against artifacts it did not produce. For each model
the corpus \code{refs/} stores the publisher's own configuration file and
tensor manifest (names and shapes), fetched reproducibly, with a provenance
audit: every file must trace to a declared source, because a reference file
without provenance is exactly the hard-coded knowledge the project exists to
remove. The checker compares the description field by field and tensor by
tensor against these artifacts. Numeric grading uses independent
implementations---the Transformers model class compiled from the publisher's
configuration, and upstream \code{llama.cpp} for exported artifacts---rather
than any code path shared with the description's compilation.

\subsection{Four support states}

\begin{table}[t]
\centering
\caption{Support states. A claim may only move upward with new evidence;
shadow (random-weight) comparisons are wiring evidence and never count as
run.}
\label{tab:states}
\begin{tabular}{@{}p{2.6cm}p{5.2cm}p{5.2cm}@{}}
\toprule
State & Meaning & Evidence \\
\midrule
Specification & config fields and tensor names/shapes checked against
publisher artifacts & per-model diff report \\
Load & real checkpoint loads with zero missing and zero unexpected tensors &
loader report from the GPU benchmark \\
Run & real-weight forward or generation matches an independent reference &
numeric comparison report \\
Artifact & exported standard artifact loads and short-generates in unmodified
upstream & round-trip report \\
\bottomrule
\end{tabular}
\end{table}

Table~\ref{tab:states} is the project's answer to a familiar marketing
failure: a model that merely \q{runs} being advertised as supported. The
support matrix records every model and feature in these four states, with
experimental entries excluded from stable counts.

\subsection{Discipline}

Three rules are enforced socially and by the check suite. \emph{Negative
examples must fail}: for every rule there is a description that violates it
and is expected to be rejected. \emph{Shadow comparisons stay wiring
evidence}: matching a random-weight shadow model bit-for-bit proves that
wiring is right and nothing else. \emph{Numbers rise only with evidence}:
the stable-coverage count is computed from the matrix by rule, not by prose.
The criteria themselves are explicitly open to audit: the repository ships a
recomputation task list asking an outsider to hand-count tensor totals and
config-field hits for three models of different sizes, and to hand-derive a
per-token key--value cache size and compare it with the checker's output. The
stated position is that the raw data come from publishers and the conclusions
from the project's own programs---only the \emph{criteria} are the thing an
outsider needs to attack.

\section{Implementation}
\label{sec:impl}

\subsection{Backends}

\begin{table}[t]
\centering
\caption{Backends and their roles. The first four are cross-checked against
each other on shared architectures; the Transformers backend executes
complete published models.}
\label{tab:backends}
\begin{tabular}{@{}p{3.2cm}p{3.2cm}p{6.2cm}@{}}
\toprule
Backend & Written in & Role \\
\midrule
NumPy & Python & dependency-light CPU execution; the first-run demo \\
PyTorch & Python & generated trainable model with a build entry point,
training steps \\
C & generated source & low-level execution path, cross-checking \\
ggml & generated source & validation backend for the GGUF path \\
Transformers & Python & complete published models: tying, vision, MTP,
quantized checkpoints \\
\bottomrule
\end{tabular}
\end{table}

Table~\ref{tab:backends} lists the execution paths. \q{Cross-checked} means a
shared architecture must produce the same forward on each backend that claims
it; the IR deliberately contains no framework concepts (no module objects, no
framework shape objects), so that this comparison is meaningful. The
Transformers backend exists because complete published models---parameter
tying, vision towers, MTP, MXFP4 and FP8 storage---should not be forced
through a lossy serial lowering; it compiles a Transformers configuration from
the declaration and uses the runtime's own operators.

MTP deserves a note: Transformers ignores a checkpoint's MTP tensors, so
\byone{} supplies the MTP computation itself (embedding and output head shared
with the trunk, one decoder, a final normalization), following the connection
formula of the vLLM implementation~\cite{kwon2023vllm}. The MTP output is
checked against an independent formula, not against \byone{}'s own reciprocal.

\subsection{Tooling}

The command line covers the loop a reviewer actually runs: specification
checking, NumPy execution, generated-PyTorch export, training steps from a
saved batch, plus check suites for the language specification, the IR entry
point, end-to-end runs, mutation testing of descriptions, and GPU smoke and
benchmark runs. The core check path is standard-library-only, so the first
tier of verification runs on a bare Python installation.

\section{Evaluation}
\label{sec:eval}

\subsection{Setup}

Measurements were collected on October 13, 2026, on a single machine with an
NVIDIA RTX PRO 6000 Blackwell Server Edition (96\,GB), PyTorch 2.8.0, CUDA
12.8, and Transformers 5.19.0. The comparison loads the same checkpoint on
both sides: the reference loads the publisher's configuration; \byone{}
compiles the configuration from the \code{.by1} declaration. Both sides use
the same Transformers model implementation and operators. Settings: BF16,
batch size 1, eager attention, input lengths 32 and 128, two warmups and five
timed runs, and greedy generation of 16 new tokens. Raw JSON reports ship with
the repository.

\subsection{Published-checkpoint parity}

\begin{table}[t]
\centering
\caption{Six published checkpoints. \q{Diff} is the maximum absolute logit
difference against the reference; every model generated identical greedy
tokens. Latency ratios are single-machine observations at length 128, not
speed claims.}
\label{tab:parity}
\begin{tabular}{@{}lrrrc@{}}
\toprule
Model & Latency ratio & Peak memory & Diff & Extra check \\
\midrule
GPT-2 & 0.990$\times$ & 0.27\,GB & 0 & 16 tokens matched \\
Qwen3-0.6B & 0.996$\times$ & 1.24\,GB & 0 & training step tested \\
Qwen2.5-7B-Instruct & 0.998$\times$ & 15.28\,GB & 0 & 16 tokens matched \\
Qwen3-8B & 0.997$\times$ & 16.43\,GB & 0 & longer contexts \\
\code{gpt-oss-20b} & 1.054$\times$ & 41.90\,GB & 0 & MXFP4 dequantized on
both sides \\
Qwen3.8-27B & 0.991$\times$ & 55.64\,GB & 0 & vision input and MTP \\
\bottomrule
\end{tabular}
\end{table}

Table~\ref{tab:parity} is the core result: across six checkpoints spanning
GPT-2 to a 27B multimodal model, including an MoE model stored in MXFP4 and
dequantized to BF16 identically on both sides, the compiled description
matches the reference \emph{exactly}. Exactness matters more than a tolerance
here: the question under test is whether the description, the configuration
conversion, the weight loading and the wiring are faithful, and for that
question any nonzero difference is a defect.

\subsection{Longer contexts and memory}

\begin{table}[t]
\centering
\caption{Long-context runs. All maximum absolute logit differences were 0.}
\label{tab:longctx}
\begin{tabular}{@{}lrrl@{}}
\toprule
Model & Input length & Peak memory & Generation \\
\midrule
Qwen2.5-7B-Instruct & 2{,}048 & 16.52\,GB & 8 tokens matched \\
Qwen2.5-7B-Instruct & 4{,}096 & 20.16\,GB & 8 tokens matched \\
Qwen3-8B & 2{,}048 & 17.85\,GB & 8 tokens matched \\
Qwen3-8B & 4{,}096 & 22.01\,GB & 8 tokens matched \\
Qwen3-8B & 8{,}192 & 38.44\,GB & 4 tokens matched \\
\bottomrule
\end{tabular}
\end{table}

\subsection{A real training step}

A single BF16 full-parameter training step on Qwen3-0.6B (real weights,
backpropagation, one AdamW update) was compared elementwise against the
official model: 596{,}049{,}920 gradients and the same number of updated
parameter values, both with maximum absolute difference 0, at about
5.98\,GB peak memory. This validates one step of loading, forward, backward
and update wiring; it is not evidence about multi-step training or
convergence.

\subsection{Vision and MTP}

The Qwen3.8-27B path uses real checkpoint weights for text, MTP and vision.
The image case feeds a synthetic $128\times128$ image through the model's
processor; maximum absolute logit difference was 0. The synthetic image
validates execution and wiring, not visual capability. MTP output matched an
independent connection formula with maximum error 0.

\subsection{Coverage}

As of the release-gate inventory, 37 model specifications are stable at the
specification state (32 additional models from the reference corpus plus six
complete models, with GPT-2 overlapping), 8 models reach the run state, and 1
(GPT-2's standard GGUF chain) reaches the artifact state. Native paths do not
execute vision towers or MTP; the audio tower is not modeled at all;
quantization has verified storage contracts (MXFP4 and FP8) and a connected
loading path, but runtime and export states are explicitly not claimed.

\subsection{Threats to validity}

The parity results reuse the reference operators, so they validate
description, conversion and loading---not kernels, and not speed. The latency
ratios are observations from one machine and short sequences; the
\code{gpt-oss-20b} ratio above 1 reflects dequantization overhead on both
sides, not a claim. Everything here is single-machine work with a small
number of timed repetitions. The criteria (thresholds, the definition of
\q{covered}, the choice of models) were written by the same project and are
the subject of the shipped audit task list. A current-source re-run of the
real-checkpoint comparison was pending at the time of writing because one
model download stalled.

\section{Lessons and Negative Results}
\label{sec:lessons}

\paragraph{Describable is not computable.} The false-positive incident
(Section~\ref{sec:problem}) is the project's most expensive lesson, and it is
a lesson about abstraction boundaries, not about test coverage: a checking
suite can only check what its contracts mention. The cure was to narrow the
IR's vocabulary so that the wrong lowering is unrepresentable, and to make
unimplemented combinations hard errors.

\paragraph{Silent fallback accumulates into fiction.} Every place where a
component guessed a default turned into a discrepancy between carriers.
Banning guessing produced better error messages and, occasionally, the
discovery that a feature had never been implemented at all.

\paragraph{Generated specifications beat maintained specifications.} The IR
document is generated from executable tables; the language specification's
examples are executed as tests. Both mechanisms cost little and remove a
whole class of drift.

\paragraph{Do not promote wiring evidence.} Random-weight shadow comparisons
are cheap and satisfying, and exactly the kind of evidence that gets
over-read. Keeping them in a separate column of the support matrix costs
nothing and keeps the claims honest.

\paragraph{Provenance.} Nearly all code, documentation and commit messages in
this project were written by an AI coding agent under human direction. This
is stated to raise scrutiny, not trust: it is precisely why the repository
pins negative-test baselines, generates its specifications, and ships
recomputation instructions for outsiders. Agent-written code makes the
argument for external graders stronger, not weaker.

\section{Related Work}

\paragraph{Model configurations and hubs.} Transformers~\cite{wolf2020transformers}
popularized the configuration-plus-weights release format that \byone{}
grades against; its configuration objects are the reference side of our
comparisons. Releases such as Gemma~\cite{gemma2024},
Qwen~\cite{qwen2025qwen25} and \code{gpt-oss}~\cite{gptoss2025} show the
diversity of architectures a single hub must express---tying,
GQA~\cite{ainslie2023gqa}, rotary position variants~\cite{su2021roformer,
peng2023yarn}, RMSNorm~\cite{zhang2019rmsnorm}, gated
FFNs~\cite{shazeer2020glu}, MoE~\cite{fedus2022switch},
SSMs~\cite{daogu2023mamba}---which is the expressiveness problem our pattern
and mechanism design addresses.

\paragraph{Exchange formats and compilers.} ONNX~\cite{onnx2019},
TVM/Relay~\cite{chen2018tvm}, MLIR~\cite{lattner2021mlir},
XLA~\cite{abadi2016tensorflow} and \code{jax} transformations~\cite{jax2018}
describe \emph{computation graphs} for compilation. \ir{} sits at a
different level: it describes \emph{architecture contracts}---which tensors
exist, how they are shared, what state is held, how layers are
scheduled---and deliberately leaves graph lowering to runtimes. This is why
its kind set is closed and tiny rather than extensible.

\paragraph{Runtimes and weight formats.} \code{llama.cpp} and
GGUF~\cite{llamacpp2023} standardize metadata and quantization for
deployment; \byone{}'s export path targets standard GGUF consumed by
unmodified upstream, and treats private extensions as a last resort. Serving
systems such as vLLM~\cite{kwon2023vllm} motivate the state contracts
(cache reuse, prefix sharing) and the MTP connection formula we validate
against. Training systems~\cite{shoeybi2019megatron} cover scale; \byone{}'s
training surface is limited to single-step verification of loading and
gradient wiring. Multi-token prediction as a training objective is studied in
DeepSeek-V3~\cite{deepseekai2024deepseekv3}.

\paragraph{Quantization formats.} MXFP4 and FP8 storage contracts follow the
OCP microscaling specification~\cite{ocp2023mx}; \byone{} checks storage
contracts and manifests rather than implementing quantizers.

\paragraph{Testing and verification of DNNs.} Whitebox testing
\cite{pei2017deepxplore} and its line attack model behavior from inputs.
\byone{} attacks the description layer instead: the artifacts under test are
configurations, manifests and contracts, and the oracle is the publisher's
own product plus an independent runtime. Weight tying itself is old
knowledge~\cite{presswolf2017using}; making it a checkable relation rather
than a code convention is the contribution.

\section{Conclusion}

\byone{} is an attempt to give the architecture of a large model a single
place to live that machines can check. The empirical claim we make is narrow
and, we think, exactly sized: for six published checkpoints, a description
compiled to a neutral IR, loaded with real weights and executed on established
runtimes, reproduces the reference numerics exactly---including a 27B model
with a vision tower and MTP, and an MoE model in MXFP4---and one real training
step matches gradient-for-gradient. The broader claim is methodological:
support should be reported in evidence states, graders should live outside
the project, unimplemented combinations should be rejected, and
specifications should be generated wherever a program can generate them. The
open work is equally clear: reaching the artifact state for the remaining
models, runtime verification of quantized paths, modeling the audio tower,
and---most importantly---external recomputation of the criteria we wrote
ourselves.

\bibliographystyle{plainnat}
\bibliography{by1-refs}

\end{document}
